\documentclass[11pt]{article}
\usepackage[a4paper,margin=1in]{geometry}
\usepackage{amsmath,amssymb,mathtools,bm}
\usepackage{enumitem,booktabs,array}
\usepackage{tcolorbox}
\usepackage{hyperref}
\usepackage[protrusion=true,expansion=false]{microtype}
\hypersetup{colorlinks=true,linkcolor=blue,urlcolor=blue,citecolor=blue}
\setlist[itemize]{topsep=2pt,itemsep=2pt}
\setlist[enumerate]{topsep=2pt,itemsep=2pt}
\newtcolorbox{keybox}{colback=blue!3!white,colframe=blue!40!black,boxrule=0.4pt,arc=1mm}
\newtcolorbox{alertbox}{colback=red!3!white,colframe=red!40!black,boxrule=0.4pt,arc=1mm}
\newtcolorbox{answerbox}{colback=green!3!white,colframe=green!40!black,boxrule=0.4pt,arc=1mm}
\newtcolorbox{drillbox}{colback=yellow!5!white,colframe=orange!60!black,boxrule=0.4pt,arc=1mm}
\newcommand{\EV}{\mathbb{E}}
\newcommand{\blank}[1]{\underline{\hspace{#1}}}

\title{W01-04: Berger \& Bouwman (2009, RFS) \\
\large ``Bank Liquidity Creation'' --- Active Recall Workbook}
\author{Banking \& Risk Management --- Exam Prep}
\date{\today}
\begin{document}\maketitle

\begin{keybox}
\textbf{Paper in one sentence.} Introduces the first comprehensive quantitative measure of bank liquidity creation based on classifying assets, liabilities, and off-balance-sheet items by liquidity, then aggregating; documents total US-bank liquidity creation of \textbf{\$2.8 trillion} in 2003, growing $\sim 80\%$ in real terms 1993--2003, concentrated in large banks and in off-balance-sheet commitments.

\textbf{Placement.} Empirical measurement paper that gives the profession its standard metric for what banks do. Complements KRS (W01-01) / GS (W01-02) theory and Ashcraft (W01-03) causal evidence.
\end{keybox}

\section*{Round 1: Complete Walk-Through}

\subsection*{1.1 The BB measure}
Every balance-sheet and off-balance-sheet item is classified as \emph{liquid}, \emph{semiliquid}, or \emph{illiquid}, and weights are assigned:
\[
\text{Liquidity Creation} = \tfrac{1}{2}\!\sum(\text{illiquid assets + liquid liabilities + illiquid OBS}) - \tfrac{1}{2}\!\sum(\text{liquid assets + illiquid liabilities + equity + liquid OBS}).
\]
Intuition: banks \emph{create liquidity} when they finance illiquid assets with liquid claims (a demand deposit backs a long-term business loan). They \emph{destroy liquidity} when they do the reverse (invest in cash with long-term equity).

\subsection*{1.2 Four measures}
BB produce four measures crossing two dimensions:
\begin{itemize}
\item Include vs.\ exclude off-balance-sheet (OBS) items.
\item Classify loans by \emph{category} (CRE, C\&I, etc.) vs.\ \emph{maturity}.
\end{itemize}
The \emph{preferred} measure is ``cat fat'' = category classification with OBS included.

\subsection*{1.3 Dataset \& headline}
\begin{itemize}
\item All US commercial banks, 1993--2003, Call Report data.
\item Aggregate bank liquidity creation: \$2.8T in 2003, up from \$1.6T (real) in 1993 --- roughly 80\% real growth.
\item Large banks (top 25 by assets) produce $>$60\% of total liquidity creation despite being $\sim$13\% of banks.
\item Off-balance-sheet items (loan commitments, standby letters of credit) account for $\sim$40\% of total liquidity creation.
\end{itemize}

\subsection*{1.4 Capital--liquidity relationship}
\begin{itemize}
\item For \emph{large} banks: liquidity creation is \emph{positively} related to capital, consistent with ``risk-absorption'' hypothesis (more capital allows more risk-taking and liquidity creation).
\item For \emph{small} banks: liquidity creation is \emph{negatively} related to capital, consistent with ``financial fragility/crowding-out'' hypothesis (Diamond--Rajan 2000/01: fragile capital structure commits bank to liquidity creation).
\item This bifurcation is the paper's most-cited cross-sectional result.
\end{itemize}

\subsection*{1.5 Identification \& caveats}
\begin{alertbox}
\begin{itemize}
\item The measure is an accounting classification, not a structural estimate. Weights of $\pm\tfrac12$ are ad hoc but defensible.
\item Capital--liquidity cross-section is correlational, not causal.
\item OBS measurement pre-SFAS 133/157 is noisy.
\item Later work (Berger--Bouwman 2017; Acharya--Mora 2015) uses the measure in crisis settings.
\end{itemize}
\end{alertbox}

\section*{Round 2: Level 1 Fill-in}
\begin{enumerate}
\item Liquidity Creation $= \tfrac12\sum(\text{illiquid assets}+\blank{2cm}+\text{illiquid OBS}) - \tfrac12\sum(\text{liquid assets}+\blank{2cm}+\text{equity}+\text{liquid OBS})$.
\item Preferred measure: ``cat fat'' = \blank{1.5cm} classification with OBS included.
\item Total US bank liquidity creation in 2003: \$\blank{1cm}T.
\item Growth 1993--2003: $\sim$\blank{1cm}\% real.
\item Sign of capital--liquidity relation: \blank{0.8cm} for large banks, \blank{0.8cm} for small banks.
\end{enumerate}

\section*{Round 3: Level 2 Prompts}
\begin{enumerate}
\item Why does a demand deposit backing a business loan \emph{create} liquidity? Write it in terms of depositor and borrower liquidity needs.
\item Contrast the ``risk-absorption'' and ``financial-fragility'' hypotheses of capital--liquidity and reconcile them with BB's size-dependent result.
\item Suppose you wanted to study how liquidity creation responded to the 2008 crisis. Sketch the empirical design using the BB measure.
\item Why is off-balance-sheet inclusion essential for the preferred measure?
\end{enumerate}

\section*{Round 4: Minimal Scaffolding}
From a blank page: (i) BB formula, (ii) four measures, (iii) 2003 aggregate and growth, (iv) bifurcated capital result, (v) two measurement concerns.

\section*{Round 5: Blank Page}
Essay: is bank liquidity creation a better measure of ``what banks do'' than total assets or loan volume? Defend or critique BB.

\vspace{6pt}
\begin{drillbox}
\textbf{Speed Drill.} (1) BB weights. (2) 2003 aggregate \$. (3) ``Cat fat'' definition. (4) Capital sign for large banks. (5) Capital sign for small banks.
\end{drillbox}

\section*{Quick Reference \& Answer Key}
\begin{answerbox}
\begin{itemize}
\item \textbf{Weights:} $+\tfrac12$ on illiquid assets, liquid liabilities, illiquid OBS; $-\tfrac12$ on liquid assets, illiquid liabilities, equity, liquid OBS.
\item \textbf{Aggregate.} \$2.8T in 2003, $\sim$80\% real growth 1993--2003.
\item \textbf{Preferred measure.} ``cat fat'': \emph{cat}egory loan classification, with \emph{fat} (OBS included).
\item \textbf{Cross-section.} Capital $+$ for large banks (risk-absorption); Capital $-$ for small banks (fragility / Diamond--Rajan).
\end{itemize}
\end{answerbox}

\begin{keybox}
\textbf{30-second exam pitch.} Berger--Bouwman (2009) give the field its workhorse measure of bank liquidity creation. The construction is straightforward: classify every asset, liability, and off-balance-sheet item as liquid, semiliquid, or illiquid, and aggregate with $\pm\tfrac12$ weights so that liquidity is created when illiquid assets (e.g., business loans) are funded by liquid claims (e.g., demand deposits). Using Call Report data 1993--2003, they document \$2.8T of US bank liquidity creation in 2003, real growth of $\sim$80\% over the decade, and a heavy concentration in large banks and off-balance-sheet items. The cross-sectional capital result is bifurcated: large banks create more liquidity when better capitalised (risk-absorption); small banks create less (Diamond--Rajan fragility). The measure is an accounting construct, but it has become the standard quantitative answer to ``what do banks do?''
\end{keybox}

\end{document}
